% 第5章：本章小结。

\section{本章小结}
\label{sec:ls-representation-summary}

本章把对象与记录对象的数分开。前四节只用线性结构。线性空间统一了可以作线性组合的对象，子空间、直和与商空间分别刻画空间内部的范围、唯一分解与被忽略的差别（第~\ref{sec:ls-vector-spaces}节）。基同时满足张成与线性无关，交换引理保证维数不依赖基的选取；选定有序基以后，对象由坐标记录，换基矩阵负责坐标换算，而近共线的基会给同一对象很大的系数（第~\ref{sec:ls-bases-dimension-coordinates}节）。线性映射在固定基下与矩阵一一对应；核与像记录丢失与保留的信息，秩零度定理把二者的维数联系起来，并决定四个基本子空间的维数与可逆性。两端独立换基得到相抵，秩是它的完整不变量；算子两端同时换基得到相似，秩、行列式与迹都只是必要条件（第~\ref{sec:ls-linear-maps}节）。线性泛函组成对偶空间，对偶映射在对偶基下的矩阵是转置（第~\ref{sec:ls-dual-spaces}节）。

后两节加入度量。有限维中一切范数等价，收敛就是逐坐标收敛，线性映射自动连续，子空间都是闭集，空间都是完备的（第~\ref{sec:ls-normed-spaces}节）。内积在长度之外给出角度：Cauchy--Schwarz不等式使夹角有意义；Gram矩阵判断线性无关，并给出坐标下的内积；正交补是由内积唯一确定的补空间，四个基本子空间由此两两配成正交直和；Gram--Schmidt正交化与QR分解造出标准正交基；Riesz表示把对偶映射变成伴随；完备的内积空间是Hilbert空间（第~\ref{sec:ls-inner-product-spaces}节）。

本章说明了记录的是什么、记录怎样随基改变。同一个算子反复作用时在哪些方向上最简单，由第~\ref{chap:linear-analysis}章的特征值与对角化回答；方程何时有解、解是否唯一、无解时怎样取最近的近似，在第~\ref{chap:linear-equations-and-optimization}章展开。贯穿例$\bm X_\varepsilon$的近共线结构将在这些章节中继续出现。
